\newpage
\thispagestyle{plain}
\begin{center}
{\Large \textbf{ABSTRACT / SYNOPSIS}}\\[1.0cm]
\end{center}

\noindent As enterprise software systems rapidly transition from passive conversational Large Language Models (LLMs) to autonomous agentic architectures capable of invoking external tools, querying databases, and transmitting emails, they introduce a critical vulnerability: \textbf{Indirect Prompt Injection} (OWASP LLM01). Because Transformer-based architectures process trusted system instructions and untrusted third-party inputs within an undifferentiated token context, malicious instructions embedded in external data can hijack agent execution, rendering the agent a ``Confused Deputy'' (OWASP LLM06). Existing defenses---including system-prompt guardrails and human-in-the-loop (HITL) approval gates---consistently fail because natural language guardrails are bypassable and human approvers are susceptible to social engineering when presented with benign-looking rationales crafted by compromised agents.

\vspace{0.4cm}
\noindent To resolve this challenge, this project introduces \textbf{AEGIS (Autonomous Entity Governance and Inspection Shield)}, an in-line security gateway and cryptographic approval proxy for autonomous AI agents. AEGIS introduces a novel security axiom: \textbf{Data Provenance Strictly Dominates Human Approval}. Grounded in a dual-lattice Information Flow Control (IFC) model uniting Denning's confidentiality lattice with Biba's integrity lattice, AEGIS mathematically proves (Theorem 1) that information originating from an untrusted source cannot flow into a privileged outbound sink, regardless of whether a human operator grants approval. 

\vspace{0.4cm}
\noindent The AEGIS gateway integrates a multi-stage inspection pipeline combining sub-millisecond heuristic filtering, ONNX-quantized neural injection detection, Presidio-based Data Loss Prevention (DLP), Open Policy Agent (OPA) rule enforcement, and SHA-256 cryptographic token binding to prevent Time-of-Check to Time-of-Use (TOCTOU) payload tampering. Empirical evaluation across 10 distinct real-world attack scenarios (E01--E10) demonstrates a 100\% detection rate with a mean gateway latency overhead of only 17.98~ms. The project passes 106 automated unit and integration tests (100\% pass rate in 2.59~s) and maintains zero database lock contention under 50-worker concurrent load via SQLite Write-Ahead Logging (WAL). A real-time Security Operations Center (SOC) dashboard provides security analysts with threat distribution analytics and causal lineage directed acyclic graphs (DAGs).

\vspace{0.8cm}
\noindent \textbf{Keywords:} AI Agent Security, Indirect Prompt Injection, Information Flow Control (IFC), Data Provenance, Confused Deputy, Cryptographic Token Binding, Dual-Lattice Security.
